\documentclass[11pt,a4paper]{article}

\usepackage[margin=2.5cm]{geometry}
\usepackage{amsmath,amssymb,amsthm}
\usepackage{enumitem}
\usepackage{hyperref}
\usepackage{tikz-cd}
\usepackage{booktabs}
\usepackage{listings}
\usepackage{xcolor}

\theoremstyle{definition}
\newtheorem{definition}{Definition}[section]
\newtheorem{example}[definition]{Example}
\theoremstyle{plain}
\newtheorem{theorem}[definition]{Theorem}
\newtheorem{lemma}[definition]{Lemma}
\newtheorem{corollary}[definition]{Corollary}
\newtheorem{proposition}[definition]{Proposition}
\theoremstyle{remark}
\newtheorem{remark}[definition]{Remark}

\newcommand{\DCont}{\mathbf{DCont}}
\newcommand{\Cat}{\mathbf{Cat}}
\newcommand{\Set}{\mathbf{Set}}
\newcommand{\Pos}{\mathrm{Pos}}
\newcommand{\sub}{\mathop{\downarrow}}
\newcommand{\shape}{\mathrm{Shape}}
\newcommand{\Ext}{\mathrm{Ext}}
\newcommand{\op}{\mathrm{op}}
\newcommand{\Ctx}{\mathbf{Ctx}}
\newcommand{\viewer}{\mathrm{viewer}}
\newcommand{\ego}{\mathrm{ego}}
\renewcommand{\root}{\mathrm{o}}

\lstdefinestyle{haskell}{
  language=Haskell,
  basicstyle=\ttfamily\small,
  keywordstyle=\color{blue!70!black}\bfseries,
  commentstyle=\color{green!50!black}\itshape,
  stringstyle=\color{red!60!black},
  showstringspaces=false,
  columns=flexible,
  breaklines=true,
  literate={->}{$\to$}2 {=>}{$\Rightarrow$}2 {<-}{$\leftarrow$}2
           {forall}{$\forall$}1 {::}{$\mathop{::}$}2
}

\lstdefinestyle{lean}{
  basicstyle=\ttfamily\small,
  keywordstyle=\color{purple!70!black}\bfseries,
  commentstyle=\color{green!50!black}\itshape,
  showstringspaces=false,
  columns=flexible,
  breaklines=true,
  morekeywords={def,theorem,structure,where,import,namespace,end,variable,
                inductive,instance,class,deriving,by,intro,exact,rfl,funext,
                show,apply,obtain,have,let,fun,Type,Prop,sorry}
}

\title{Social Networks as Directed Container Comonads:\\
A Categorical Framework with Controlled Sheaf Failure}
\author{Robin Langer \and Neil Ghani}
\date{July 2026}

\begin{document}
\maketitle

\begin{abstract}
We present a categorical framework for social networks built on three
interlocking structures: (1)~\emph{directed containers}
(Ahman--Chapman--Uustalu) model the social graph, where the comonad's
\texttt{extract} produces a user's feed and \texttt{extend} applies
feed algorithms across all perspectives; (2)~\emph{presheaves} on a
category of viewing contexts model permissions, with contravariance
capturing the principle that less privilege sees less content;
(3)~the \emph{sheaf condition} classifies content by consistency
guarantee---public posts form a sheaf (all viewers agree), while
blocked content and personalised feeds deliberately fail the gluing
axiom. The framework is implemented in Haskell (Servant, type-safe
API) with invariants verified in sorry-free Lean~4, reusing the
container and comonad library from the \emph{trust-boundaries}
project. The Haskell types serve double duty: they are the runtime
data model \emph{and} the template for Lean theorem statements, with
the Curry--Howard correspondence making the translation nearly
mechanical.
\end{abstract}

\tableofcontents

%% ===================================================================
\section{Introduction}
%% ===================================================================

Social networks are widely acknowledged to be hard to build correctly.
The difficulty is not any single problem but the \emph{product} of
several interacting concerns: graph traversal, access control,
consistency, and ranking. We argue that these concerns decompose
cleanly along categorical lines:

\begin{enumerate}[nosep]
  \item The \textbf{feed problem}---computing a personalised view of
    the graph---is an instance of the comonadic \texttt{extend}
    operation on a directed container.
  \item The \textbf{permission problem}---who sees what---is the
    functoriality condition of a presheaf on viewing contexts.
  \item The \textbf{consistency problem}---where do different viewers
    agree?---is the sheaf condition, and a well-designed social network
    controls \emph{exactly where} this condition holds and where it
    fails.
\end{enumerate}

This paper develops the mathematics, provides a Haskell implementation
that mirrors the categorical structure, and identifies the proof
obligations that we discharge in Lean~4.

\paragraph{Relation to prior work.}
The directed container library (containers, directed containers,
comonad laws, coKleisli composition, and strength) was developed and
formalised in the \emph{trust-boundaries} project~\cite{trust}, where
proof registries for epistemically honest AI agents were modelled as
directed container comonads. We reuse that library wholesale and
provide a new instantiation for social graphs.

%% ===================================================================
\section{Background: containers and directed containers}
\label{sec:background}
%% ===================================================================

We recall the definitions from Abbott--Altenkirch--Ghani~\cite{aag}
and Ahman--Chapman--Uustalu~\cite{acu}, following the Lean
formalisation in the trust-boundaries library.

\begin{definition}[Container]
A \emph{container} is a pair $C = (S, P)$ where $S$ is a type of
\emph{shapes} and $P : S \to \mathrm{Type}$ is a family of
\emph{positions}. Its \emph{extension} is the polynomial functor
\[
  \llbracket C \rrbracket(X) = \sum_{s : S} (P(s) \to X).
\]
\end{definition}

Container morphisms $(f, f^\sharp) : C \to D$ consist of a shape map
$f : S_C \to S_D$ and a contravariant position map
$f^\sharp_s : P_D(f\,s) \to P_C(s)$. The representation theorem
establishes: container morphisms $C \to D$ are in bijection with
natural transformations
$\llbracket C \rrbracket \Rightarrow \llbracket D \rrbracket$.

\begin{definition}[Directed container]\label{def:dcont}
A \emph{directed container} extends a container $(S, P)$ with three
operations:
\begin{itemize}[nosep]
  \item a \emph{root} $\root : \prod_{s:S} P(s)$;
  \item a \emph{sub-shape} ${\sub} : \prod_{s:S} P(s) \to S$;
  \item a \emph{shift} $\oplus : \prod_{s:S} \prod_{p:P(s)} P(s \sub p) \to P(s)$;
\end{itemize}
subject to five laws:
\begin{enumerate}[label=\textbf{D\arabic*},nosep]
  \item $s \sub \root(s) = s$
  \item $\root(s) \oplus q = q$ \quad(transported along D1)
  \item $p \oplus \root(s \sub p) = p$
  \item $s \sub (p \oplus q) = (s \sub p) \sub q$
  \item $(p \oplus q) \oplus r = p \oplus (q \oplus r)$ \quad(transported along D4)
\end{enumerate}
\end{definition}

\begin{theorem}[Comonad structure {\cite{acu}}]\label{thm:comonad}
A directed container $(S, P, \root, {\sub}, \oplus)$ induces a comonad
on $\llbracket C \rrbracket$ with:
\begin{align*}
  \varepsilon\,\langle s, v \rangle &= v(\root(s)) \\
  \delta\,\langle s, v \rangle &= \langle s,\; \lambda p.\, \langle s \sub p,\; \lambda q.\, v(p \oplus q) \rangle \rangle
\end{align*}
The three comonad laws follow from D1--D5:
\begin{itemize}[nosep]
  \item Left counit $\varepsilon \circ \delta = \mathrm{id}$ from D1, D2.
  \item Right counit $\llbracket C \rrbracket(\varepsilon) \circ \delta = \mathrm{id}$ from D3.
  \item Coassociativity $\llbracket C \rrbracket(\delta) \circ \delta = \delta \circ \delta$ from D4, D5.
\end{itemize}
All five laws and all three comonad laws are formalised sorry-free in
Lean~4 (files \texttt{Directed.lean}, \texttt{CoKleisli.lean} in the
trust-boundaries library).
\end{theorem}

%% ===================================================================
\section{The social graph as a directed container}
\label{sec:social-dcont}
%% ===================================================================

We now instantiate the abstract framework for social networks.

\begin{definition}[Social graph container]\label{def:social-graph}
Let $\mathcal{G}$ be the directed container defined by:
\begin{itemize}[nosep]
  \item $S = $ the type of \emph{social graphs}: a set of users $U$
    together with a follow relation $F \subseteq U \times U$ and a
    distinguished viewer $v_0 \in U$.
  \item $P(s) = U_s$, the set of users in graph $s$.
  \item $\root(s) = v_0$, the logged-in viewer.
  \item $s \sub p = \ego(s, p)$, the \emph{ego network} of user $p$
    in graph $s$: the subgraph induced by $p$ and all users $p$ follows.
  \item $p \oplus q$: given $p \in P(s)$ and $q \in P(s \sub p)$,
    the user $q$ (which is also a user in $s$, since $P(s \sub p) \subseteq P(s)$).
\end{itemize}
\end{definition}

\begin{proposition}[D1--D5 for social graphs]\label{prop:social-d1d5}
The social graph container satisfies the five directed container laws:
\begin{enumerate}[label=\textbf{D\arabic*},nosep]
  \item $\ego(s, v_0) = s$: the ego network of the viewer is the
    entire visible graph. (Holds when the viewer can see everyone they
    follow---the public-posts assumption.)
  \item $v_0 \oplus q = q$: navigating from the viewer to position $q$
    in the viewer's ego network yields $q$. (Immediate from the
    definition of $\oplus$ as inclusion.)
  \item $p \oplus \root(\ego(s,p)) = p \oplus p = p$: navigating from $p$
    to $p$'s own root position returns to $p$. (The root of
    $\ego(s,p)$ is $p$ itself.)
  \item $\ego(s, p \oplus q) = \ego(\ego(s, p), q)$: the ego network of
    $q$ in the full graph equals the ego network of $q$ within $p$'s
    ego network. (Holds when follow relations are global---$q$'s
    follows don't depend on who's looking.)
  \item Shift associativity: $(p \oplus q) \oplus r = p \oplus (q \oplus r)$.
    (Both sides equal $r$, since $\oplus$ is inclusion.)
\end{enumerate}
\end{proposition}

\begin{remark}
D4 is the most substantive law for social networks. It says that the
ego network is \emph{context-independent}: what user $q$ can see
doesn't depend on which larger graph you computed $q$'s ego network
within. This holds for public, asymmetric-follow networks (like early
Twitter). It \emph{fails} when privacy settings make visibility
context-dependent---which is exactly where the presheaf layer
(Section~\ref{sec:presheaf}) takes over.
\end{remark}

\begin{corollary}[Social feed as comonadic extract]
By Theorem~\ref{thm:comonad}, the social graph comonad has:
\begin{align*}
  \varepsilon\,\langle s, v \rangle &= v(v_0) && \text{(``my content''---the viewer's own posts)} \\
  \delta\,\langle s, v \rangle &= \langle s,\; \lambda p.\, \langle \ego(s,p),\; \lambda q.\, v(q) \rangle \rangle && \text{(``feed of feeds'')}
\end{align*}
The \emph{feed algorithm} is a coKleisli morphism
$f : \llbracket \mathcal{G} \rrbracket(X) \to Y$---a function that
inspects a local neighbourhood and produces a summary. Applying it
everywhere via $\mathrm{extend}\,f$ produces each user's personalised
feed.
\end{corollary}

\subsection{Haskell encoding}

The Haskell types mirror the Lean definitions:

\begin{lstlisting}[style=haskell]
-- Domain types
newtype UserId = UserId Int
data User = User { userId :: UserId, handle :: Text, displayName :: Text }
data Post = Post { postId :: PostId, author :: UserId, body :: Text, createdAt :: UTCTime }

-- The social graph
data SocialGraph = SocialGraph
  { users   :: Set UserId
  , follows :: Set (UserId, UserId)  -- (follower, followed)
  , viewer  :: UserId
  }

-- Extension functor
data Ext x = Ext SocialGraph (UserId -> x)
  deriving Functor

-- Directed container operations
root :: SocialGraph -> UserId
root = viewer

egoNetwork :: SocialGraph -> UserId -> SocialGraph
egoNetwork g u = SocialGraph
  { users   = Set.insert u (followedBy u g)
  , follows = Set.filter (\(a,_) -> a == u) (follows g)
  , viewer  = u
  }

shift :: SocialGraph -> UserId -> UserId -> UserId
shift _ _ q = q  -- inclusion: q is already a user in the parent graph

-- Comonad operations (derived, not hand-written)
extract :: Ext x -> x
extract (Ext g v) = v (root g)

duplicate :: Ext x -> Ext (Ext x)
duplicate (Ext g v) = Ext g (\p -> Ext (egoNetwork g p) (\q -> v (shift g p q)))

extend :: (Ext x -> y) -> Ext x -> Ext y
extend f (Ext g v) = Ext g (\p -> f (Ext (egoNetwork g p) (\q -> v (shift g p q))))
\end{lstlisting}

\subsection{Lean proof obligations}

The abstract comonad laws are already proved in the trust-boundaries
library. The only new proof work is instantiating D1--D5 for the
social graph:

\begin{lstlisting}[style=lean]
-- New file: SocialGraph.lean
import Containers.Directed

def SocialGraph : DirectedContainer where
  Shape := GraphShape           -- users + follow relation + viewer
  Pos   := fun g => g.Users
  root  := fun g => g.viewer
  sub   := fun g u => egoNetwork g u
  shift := fun g _ q => q       -- inclusion
  d1    := by intro g; ext; simp [egoNetwork]  -- ego of viewer = full graph
  d2    := by intro g q; rfl                    -- root shift = id
  d3    := by intro g p; rfl                    -- shift to sub-root = id
  d4    := by intro g p q; ext; simp [egoNetwork]  -- ego nesting
  d5    := by intro g p q r; rfl                -- shift associativity (trivial)
\end{lstlisting}

%% ===================================================================
\section{Presheaves: the permission layer}
\label{sec:presheaf}
%% ===================================================================

The directed container gives us the feed. But a real social network
must answer: \emph{which content is visible to which viewer?} This is
the domain of presheaves.

\begin{definition}[Category of contexts]\label{def:ctx}
Let $\Ctx$ be the category where:
\begin{itemize}[nosep]
  \item Objects are \emph{viewing contexts}: triples $(u, t, \pi)$
    where $u$ is a user, $t$ is a timestamp, and $\pi$ is a
    permission set (e.g., \{public, follower, admin\}).
  \item A morphism $c_1 \to c_2$ exists when $c_2$ is a
    \emph{restriction} of $c_1$: same or later time, same or fewer
    permissions, same or more restricted viewer.
\end{itemize}
Composition is transitive restriction. Identity is the trivial
restriction.
\end{definition}

\begin{definition}[Content presheaf]\label{def:presheaf}
A \emph{content presheaf} is a functor $F : \Ctx^\op \to \Set$.
For each context $c$, $F(c)$ is the set of content visible in that
context. Contravariance means: if $c_1 \to c_2$ is a restriction
(forgetting privilege), then $F(c_2 \to c_1) : F(c_1) \to F(c_2)$
maps the larger visible set into the smaller one.
\end{definition}

\begin{remark}[Privacy bugs as functoriality failures]
A privacy bug is a violation of functoriality: a case where
restricting privilege does not correctly restrict visible content. For
example, if a user blocks another but the blocked user can still see
their posts through a cached API response, the restriction morphism
has failed to propagate. The presheaf framework makes this precise:
every privacy invariant is a functoriality condition.
\end{remark}

\subsection{Haskell encoding}

\begin{lstlisting}[style=haskell]
-- Viewing context
data Context = Context
  { ctxViewer :: UserId
  , ctxTime   :: UTCTime
  , ctxPerms  :: PermissionSet
  }

-- Permission levels (ordered by privilege)
data Permission = Public | Follower | Mutual | Admin
  deriving (Eq, Ord)

-- Context morphism: evidence that c2 is a restriction of c1
data Restriction = Restriction
  { from :: Context
  , to   :: Context
  }

-- Content presheaf: contravariant functor from contexts to sets of content
class ContentPresheaf f where
  view     :: Context -> f content -> Set content
  restrict :: Restriction -> Set content -> Set content
  -- Laws (enforced in Lean):
  --   restrict id = id
  --   restrict (g . f) = restrict f . restrict g
\end{lstlisting}

\subsection{Lean formalisation}

\begin{lstlisting}[style=lean]
structure Context where
  viewer : UserId
  time   : Nat        -- discrete timestamp
  perms  : Finset Permission

-- Restriction: c1 has at least as much privilege as c2
def Restricts (c1 c2 : Context) : Prop :=
  c2.perms <=  c1.perms /\ c1.time <=  c2.time

structure ContentPresheaf where
  Content : Type
  view : Context -> Finset Content
  map_restrict : forall c1 c2, Restricts c1 c2 -> view c1 >=  view c2
  -- Contravariance: more privilege shows more content
\end{lstlisting}

%% ===================================================================
\section{The sheaf condition and its failure}
\label{sec:sheaf}
%% ===================================================================

A presheaf $F : \Ctx^\op \to \Set$ satisfies the \emph{sheaf
condition} with respect to a coverage $\{c_i \to c\}$ when compatible
local sections glue to a unique global section. In social network
terms: if every user's local view of a piece of content is
``the same,'' there is a single global truth.

\begin{definition}[Compatibility]
Given a cover $\{c_i\}_{i \in I}$ of a context $c$, a family of
local sections $\{s_i \in F(c_i)\}$ is \emph{compatible} when for
every pair $i, j$ with overlap $c_{ij} = c_i \times_c c_j$:
\[
  F(c_i \leftarrow c_{ij})(s_i) = F(c_j \leftarrow c_{ij})(s_j).
\]
\end{definition}

\subsection{Features that preserve the sheaf}

The minimal social network (Section~\ref{sec:minimal}) has a
\emph{non-constant} content presheaf: the feed query restricts
visibility to posts by users the viewer follows, so different
viewers see different sets of posts. Nevertheless, the sheaf
condition holds: for any post $p$, all viewers who can see $p$
agree on its content (text, author, timestamp), and viewers who
cannot see $p$ simply have no local section. There are no
conflicts.

Several natural features preserve this property. We classify them
by which categorical layer they modify:

\begin{center}
\begin{tabular}{@{}llll@{}}
\toprule
\textbf{Feature} & \textbf{What changes} & \textbf{Categorical layer} & \textbf{Sheaf?} \\
\midrule
Follow/unfollow & Which graph edges I traverse & Shape (directed container) & Yes \\
Mute & Which coKleisli morphism I run & CoKleisli morphism & Yes \\
Protected account & Who can create a follow edge & Gate on shape modification & Yes \\
\bottomrule
\end{tabular}
\end{center}

\begin{remark}[Muting is coKleisli, not presheaf]
Muting---hiding a followed user's posts from your feed without
unfollowing---does not alter the content presheaf. The muted user's
posts still \emph{exist} for the muter (accessible via direct
navigation). Muting only changes the feed algorithm, i.e., the
coKleisli morphism $f : \llbracket \mathcal{G} \rrbracket(X) \to Y$
that the viewer applies. The social graph (shape) retains the follow
edge, so the muted user sees no change in their follower count. This
is reader-side control: it parameterises the coKleisli morphism
without touching the presheaf.
\end{remark}

\begin{remark}[Protected accounts preserve the sheaf]
A protected account requires follow approval, but does not change the
visibility model. In both the unprotected and protected cases, the
invariant is: followers see all posts, non-followers see none. The
only difference is the mechanism by which an edge enters the follow
relation (unilateral vs.\ approval). The presheaf fibres are
unchanged; only the \emph{gate} on shape modification changes.
\end{remark}

\subsection{The minimal sheaf-breaking feature: audience selection}

The sheaf holds as long as the following invariant is maintained:

\medskip
\noindent\textbf{Uniform follower visibility:} \emph{If you follow
user $u$, you see all of $u$'s posts.}
\medskip

\noindent The sheaf breaks at the precise moment this invariant
weakens to: \emph{If you follow user $u$, you see \textbf{some} of
$u$'s posts.} The minimal feature that causes this is
\textbf{audience selection}: the author chooses, per post, which
subset of their followers can see it.

\begin{definition}[Audience-selected post]\label{def:audience}
An \emph{audience-selected post} extends the post type with an
explicit audience:
\[
  \mathrm{Post} = (\mathrm{author},\; \mathrm{body},\; \mathrm{timestamp},\; \mathrm{audience} \subseteq \mathrm{Followers}(\mathrm{author}))
\]
When $\mathrm{audience} = \mathrm{Followers}(\mathrm{author})$, the
post is visible to all followers (the sheaf case). When
$\mathrm{audience} \subsetneq \mathrm{Followers}(\mathrm{author})$,
some followers are excluded.
\end{definition}

\begin{theorem}[Audience selection violates the sheaf condition]\label{thm:audience-antisheaf}
Let Alice have followers Bob and Carol. Let $p$ be a post by Alice
with $\mathrm{audience}(p) = \{\text{Carol}\}$. Then:
\begin{itemize}[nosep]
  \item $p \in F(c_{\text{Carol}})$ --- Carol sees the post;
  \item $p \notin F(c_{\text{Bob}})$ --- Bob does not see the post;
\end{itemize}
yet Bob and Carol are both followers of Alice. No global section over
the cover $\{c_{\text{Bob}}, c_{\text{Carol}}\}$ is compatible with
both local views.
\end{theorem}

\begin{proof}
A global section must either include $p$ or exclude it. If it
includes $p$, it contradicts $p \notin F(c_{\text{Bob}})$. If it
excludes $p$, it contradicts $p \in F(c_{\text{Carol}})$. The gluing
axiom fails.
\end{proof}

\begin{remark}[One field, one table]
In the data model, audience selection adds one junction table:
\begin{center}
\texttt{post\_audience(post\_id, user\_id)}
\end{center}
and changes the feed query to:
\begin{lstlisting}[language=SQL,basicstyle=\ttfamily\small]
SELECT p.* FROM posts p
  JOIN follows f ON p.author_id = f.followed_id
  LEFT JOIN post_audience pa ON pa.post_id = p.id
  WHERE f.follower_id = :me
    AND (pa.post_id IS NULL          -- no audience restriction
         OR pa.user_id = :me)        -- or I am in the audience
  ORDER BY p.created_at DESC LIMIT 50;
\end{lstlisting}
The sheaf breaks with one additional field on the post (or
equivalently, one junction table). This is the exact cost of leaving
the sheaf.
\end{remark}

\subsection{Consistency classification}

\begin{definition}[Consistency levels]\label{def:consistency}
We classify content by sheaf behaviour:
\begin{center}
\begin{tabular}{@{}lll@{}}
\toprule
\textbf{Level} & \textbf{Meaning} & \textbf{Examples} \\
\midrule
\textsc{Sheaf} & All viewers in cover agree & Public posts, profiles, follows \\
\textsc{EventualSheaf} & Viewers converge over time & Like counts, follower counts \\
\textsc{Presheaf} & Deliberate within-follower disagreement & Audience-selected posts \\
\bottomrule
\end{tabular}
\end{center}
\end{definition}

\begin{theorem}[Public posts form a sheaf]\label{thm:public-sheaf}
Let $F$ be the content presheaf restricted to posts with
$\mathrm{audience} = \mathrm{Followers}(\mathrm{author})$ (i.e., no
audience restriction). Then $F$ satisfies the sheaf condition: for
any cover of follower contexts, all local sections are compatible,
and a unique global section exists.
\end{theorem}

\begin{proof}
For any post $p$ with unrestricted audience, every follower of the
author sees the same content (text, author, timestamp). The
restriction maps act as identity on content. Compatibility is
automatic, and the global section is $p$ itself.
\end{proof}

\subsection{Repairing the sheaf: local systems and audience fibres}

When audience selection breaks the sheaf, we do not discard the
presheaf framework. Instead, we \emph{refine} it.

\begin{definition}[Audience fibre]\label{def:audience-fibre}
For a post $p$ with audience $A \subseteq \mathrm{Followers}(\mathrm{author}(p))$,
the \emph{audience fibre} is the sub-presheaf $F_A$ obtained by
restricting the context category to viewers in $A$:
\[
  F_A(c) = \begin{cases}
    F(c) & \text{if } \viewer(c) \in A \\
    F(c) \setminus \{p\} & \text{otherwise}
  \end{cases}
\]
\end{definition}

\begin{proposition}[Audience fibres are sheaves]\label{prop:fibre-sheaf}
Within an audience fibre $F_A$, the sheaf condition holds: all
viewers in $A$ agree on the content of $p$, and all viewers outside
$A$ agree that $p$ does not exist. The failure is only
\emph{between} fibres.
\end{proposition}

\begin{remark}[Local system structure]
The audience-selected presheaf decomposes as a \emph{local system}:
a collection of sheaves (one per audience group) glued along the
subsets where audiences overlap. The obstruction to a global sheaf
is precisely the audience partition---different subsets of followers
live in different fibres. This decomposition gives a natural
architecture for implementation: cache and replicate within an
audience fibre (safe---it's a sheaf), but never across fibres
(unsafe---the sheaf condition fails).
\end{remark}

\subsection{Haskell encoding}

\begin{lstlisting}[style=haskell]
-- A post with optional audience restriction
data Post = Post
  { postId      :: PostId
  , postAuthor  :: UserId
  , postBody    :: Text
  , postCreated :: UTCTime
  , postAudience :: Maybe (Set UserId)  -- Nothing = all followers
  }

-- Consistency levels as a type-level tag
data ConsistencyLevel = Sheaf | EventualSheaf | Presheaf

-- Tag each content type with its consistency guarantee
type family Consistency (a :: Type) :: ConsistencyLevel where
  Consistency PublicPost         = 'Sheaf          -- audience = Nothing
  Consistency UserProfile        = 'Sheaf
  Consistency LikeCount          = 'EventualSheaf
  Consistency AudiencePost       = 'Presheaf       -- audience = Just subset

-- The type system prevents treating Presheaf content as Sheaf content:
-- you cannot cache or share AudiencePost across viewers without a
-- type error.
\end{lstlisting}

%% ===================================================================
\section{Interaction: comonad meets presheaf}
\label{sec:interaction}
%% ===================================================================

The directed container comonad (Section~\ref{sec:social-dcont}) and the
content presheaf (Section~\ref{sec:presheaf}) interact through the
feed computation. The feed is the comonadic \texttt{extend} applied to
a presheaf-filtered content stream.

\begin{definition}[Filtered feed]
Given a feed algorithm $f : \llbracket \mathcal{G} \rrbracket(\mathrm{Post}) \to [\mathrm{Post}]$
and a content presheaf $F$, the \emph{filtered feed} for context $c$ is:
\[
  \mathrm{feed}(c) = f\bigl(\langle \ego(s, \viewer(c)),\; \lambda p.\, F(c)(p) \rangle\bigr)
\]
where $F(c)(p)$ returns the posts by user $p$ visible in context $c$.
\end{definition}

\begin{remark}[Colimits don't commute with filtering]
A na\"ive implementation might compute the feed (a colimit of
content sources) and then filter by permissions. But $\mathrm{extend}$
does not commute with presheaf restriction in general: filtering
after ranking may produce different results than filtering before
ranking, since the ranking function sees a different neighbourhood.
This non-commutativity is a source of real bugs in social network
implementations (e.g., a ``top post'' being filtered out after ranking,
leaving a gap).
\end{remark}

%% ===================================================================
\section{The minimal social network}
\label{sec:minimal}
%% ===================================================================

We identify the \emph{irreducible core} of a social network: the
minimum feature set that exhibits network effects.

\begin{definition}[Minimal social network]
A minimal social network consists of:
\begin{enumerate}[nosep]
  \item \textbf{Identity}: a user type with handle and display name.
  \item \textbf{Asymmetric connection}: a follow relation $F \subseteq U \times U$
    (directed, no mutual consent required).
  \item \textbf{Broadcast}: a post type (text, author, timestamp).
  \item \textbf{Chronological feed}: the comonadic extract applied to
    a reverse-chronological ordering.
\end{enumerate}
\end{definition}

The data model is three tables:

\begin{center}
\begin{tabular}{@{}ll@{}}
\toprule
\textbf{Table} & \textbf{Columns} \\
\midrule
\texttt{users} & \texttt{id}, \texttt{handle}, \texttt{display\_name}, \texttt{bio}, \texttt{created\_at} \\
\texttt{follows} & \texttt{follower\_id}, \texttt{followed\_id}, \texttt{created\_at} \\
\texttt{posts} & \texttt{id}, \texttt{author\_id}, \texttt{body} (text, $\le 280$ chars), \texttt{created\_at} \\
\bottomrule
\end{tabular}
\end{center}

The feed query is a single join:

\begin{lstlisting}[language=SQL,basicstyle=\ttfamily\small]
SELECT posts.* FROM posts
  JOIN follows ON posts.author_id = follows.followed_id
  WHERE follows.follower_id = :me
  ORDER BY posts.created_at DESC
  LIMIT 50;
\end{lstlisting}

\begin{remark}[Categorical simplification]
The minimal social network has a non-constant content presheaf:
different viewers see different sets of posts (those by users they
follow). Nevertheless, the sheaf condition holds because the
\emph{uniform follower visibility} invariant is satisfied---every
follower of a given author sees all of that author's posts, and all
viewers who see a post agree on its content. The presheaf is
non-constant but is a sheaf, and the comonad (directed container) is
the entire story. Complexity enters only when audience selection
(Section~\ref{sec:sheaf}) breaks the uniform visibility invariant.
\end{remark}

%% ===================================================================
\section{Implementation architecture}
\label{sec:architecture}
%% ===================================================================

\begin{center}
\begin{tikzcd}[column sep=large, row sep=large]
  \text{Lean (proofs)} \arrow[d, "\text{types} \cong" swap]
  & \text{Containers.Directed} \arrow[r] \arrow[d, dashed]
  & \text{SocialGraph instance} \arrow[d, dashed] \\
  \text{Haskell (runtime)}
  & \text{DirectedContainer class} \arrow[r]
  & \text{Servant API}
\end{tikzcd}
\end{center}

The implementation proceeds in layers:

\begin{enumerate}[nosep]
  \item \textbf{Haskell domain types}: algebraic data types for
    \texttt{User}, \texttt{Post}, \texttt{SocialGraph}, \texttt{Context}.
  \item \textbf{Haskell directed container}: \texttt{root},
    \texttt{egoNetwork}, \texttt{shift} implementing the comonad.
  \item \textbf{Lean instantiation}: prove D1--D5 for these
    definitions, importing the trust-boundaries library. Comonad laws
    follow for free.
  \item \textbf{Servant API}: six type-safe HTTP routes wiring the
    comonad into a web service.
  \item \textbf{Presheaf/sheaf layer}: added when the network supports
    non-public content.
\end{enumerate}

The Haskell--Lean correspondence is nearly mechanical:

\begin{center}
\begin{tabular}{@{}ll@{}}
\toprule
\textbf{Haskell} & \textbf{Lean} \\
\midrule
\texttt{data T = A | B | C} & \texttt{inductive T | A | B | C} \\
\texttt{class C a where} & \texttt{class C (a : Type) where} \\
\texttt{instance C T where} & \texttt{instance : C T where} \\
\texttt{newtype F a = F (a -> b)} & \texttt{def F (a : Type) := a -> b} \\
type constraint & \texttt{theorem} / \texttt{Prop} \\
\bottomrule
\end{tabular}
\end{center}

%% ===================================================================
\section{Future directions}
\label{sec:future}
%% ===================================================================

\begin{enumerate}
  \item \textbf{Zappa--Sz\'ep products}: the trust-boundaries library
    includes Zappa--Sz\'ep products of directed containers
    (\texttt{ZappaSzep.lean}). In the social network setting, this
    could model the interaction between two independent graph
    structures (e.g., follow graph and group membership graph) whose
    directed containers compose non-trivially.

  \item \textbf{CoKleisli morphisms as feed algorithms}: different feed
    algorithms (chronological, ranked, ``following'' vs ``for you'')
    are different coKleisli morphisms. The coKleisli category gives a
    principled space of feed algorithms, with composition corresponding
    to pipeline stages.

  \item \textbf{Controlled sheaf failure as a design language}: rather
    than ad-hoc privacy rules, express each new feature as a choice of
    consistency level (Definition~\ref{def:consistency}). The Lean
    proof obligation forces the implementer to verify that the chosen
    level is achievable.

  \item \textbf{Cofunctors for cross-network interaction}: container
    morphisms between different social network instances
    (federation, cross-posting) are cofunctors between directed
    containers, with the trust-boundaries transport theorem
    guaranteeing that trust levels translate safely.
\end{enumerate}

%% ===================================================================
\section*{Acknowledgements}
%% ===================================================================

The directed container library and its Lean formalisation were
developed in the trust-boundaries project. The categorical
perspective on social networks emerged from conversations between
Robin Langer and Neil Ghani at Kodamai.

\begin{thebibliography}{9}

\bibitem{aag}
M.~Abbott, T.~Altenkirch, and N.~Ghani.
\newblock Containers: constructing strictly positive types.
\newblock \emph{Theoretical Computer Science}, 342(1):3--27, 2005.

\bibitem{acu}
D.~Ahman, J.~Chapman, and T.~Uustalu.
\newblock When is a container a comonad?
\newblock \emph{Foundations of Software Science and Computation Structures (FoSSaCS)}, 2012.

\bibitem{trust}
R.~Langer.
\newblock Trust boundaries: epistemically honest AI agents via directed
container comonads.
\newblock Technical report, 2026.

\bibitem{gkt}
T.~J.~Green, G.~Karvounarakis, and V.~Tannen.
\newblock Provenance semirings.
\newblock \emph{PODS}, 2007.

\end{thebibliography}

\end{document}
